% !PRINTED 511
\weberpage{511}

Eine Einheit ist ein ganzes Functional, dessen reciprokes gleichfalls ganz ist, und dieses reciproke Functional ist selbst eine Einheit. Durch eine Einheit ist jedes beliebige ganze Functional theilbar. Ueberhaupt gilt der Satz:

9. \gesperrt{Das Product und der Quotient zweier Einheiten sind wieder Einheiten.}

Denn sind $\varepsilon_1, \varepsilon_2$ zwei Einheiten, so sind $\varepsilon_1\varepsilon_2$ und $1:\varepsilon_1\varepsilon_2$ ganze Functionale, also $\varepsilon_1\varepsilon_2$ eine Einheit, und ebenso sind $\varepsilon_1:\varepsilon_2$ und $\varepsilon_2:\varepsilon_1$ ganz.

10. \gesperrt{Ist ein ganzes Functional $\mu$ theilbar durch ein anderes, $\alpha$, so ist jedes mit $\mu$ associirte Functional $\mu'$ auch durch jedes mit $\alpha$ associirte Functional $\alpha'$ theilbar.}

Denn wenn $\mu:\alpha$ ganz und $\varepsilon, \varepsilon_1$ Einheiten sind, so ist auch $\mu\varepsilon:\alpha\varepsilon_1$ ganz.

11. \gesperrt{Ist $\alpha$ associirt mit $\beta$ und mit $\gamma$, so sind auch $\beta$ und $\gamma$ unter einander associirt.}

Denn ist $\beta=\alpha\varepsilon$, $\gamma=\alpha\varepsilon_1$, so ist $\beta=\gamma\varepsilon:\varepsilon_1$, und $\varepsilon:\varepsilon_1$ ist nach 8. eine Einheit.

Ist $\alpha$ theilbar durch $\beta$, so ist die absolute Norm von $\alpha$ theilbar durch die absolute Norm von $\beta$, denn aus $\alpha=\beta\gamma$ folgt nach \S.\,136, 5.:

\weq{1}{N_a(\alpha) = N_a(\beta)\, N_a(\gamma).}

Daraus ergiebt sich weiter, dass die absolute Norm einer Einheit $=1$ sein muss. Es gilt aber auch das Umgekehrte:

12. \gesperrt{Ein ganzes Functional, dessen absolute Norm $=1$ ist, ist eine Einheit.}

Denn ist $\varepsilon$ ein ganzes Functional mit der absoluten Norm 1, so ist $N(\varepsilon)$ ein ganzes rationales Functional mit dem absoluten Werthe 1, und daher ist auch $1:N(\varepsilon)$ ein ganzes Functional. Setzen wir dann
\[
N(\varepsilon) = \varepsilon\varepsilon',
\]
so ist $\varepsilon'$ als Product von ganzen Functionalen (den Conjugirten zu $\varepsilon$) selbst ganz, und folglich ist auch
\[
\frac{1}{\varepsilon} = \frac{\varepsilon'}{N(\varepsilon)}
\]
ein ganzes Functional, also $\varepsilon$ eine Einheit.
